\section*{Теория}
\addcontentsline{toc}{section}{Теория}

\subsection*{1. Матричные разложения: QR и SVD}

\begin{theorem}[QR-разложение]
Любую невырожденную вещественную квадратную матрицу $A$ можно представить в виде произведения
\[
A = QR,
\]
где $Q$ — ортогональная матрица ($Q^T Q = E$), а $R$ — верхнетреугольная матрица с положительными элементами на главной диагонали.
\end{theorem}

\noindent\textbf{Практический алгоритм QR-разложения:}
\begin{enumerate}
    \item Столбцы матрицы $A$ рассматриваем как базис $a_1, \dots, a_n$.
    \item Применяем процесс ортогонализации Грама-Шмидта, получая ортогональную систему $b_1, \dots, b_n$, а затем нормируем её, получая ортонормированный базис $e_1, \dots, e_n$.
    \item Матрица $Q$ составляется из столбцов $e_1, \dots, e_n$.
    \item Матрица $R$ вычисляется по формуле $R = Q^T A$ (или заполняется коэффициентами разложения векторов $a_i$ по базису $e_j$).
\end{enumerate}

\begin{theorem}[Сингулярное разложение / SVD]
Любую вещественную матрицу $A$ размера $m \times n$ можно представить в виде
\[
A = U \Sigma V^T,
\]
где $U$ — ортогональная матрица размера $m \times m$, $V$ — ортогональная матрица размера $n \times n$, а $\Sigma$ — матрица размера $m \times n$, у которой на главной диагонали стоят неотрицательные числа $\sigma_1 \ge \sigma_2 \ge \dots \ge \sigma_r > 0$ (сингулярные числа), а остальные элементы равны нулю.
\end{theorem}

\noindent\textbf{Практический алгоритм нахождения SVD:}
\begin{enumerate}
    \item Вычисляем симметрическую матрицу $A^T A$.
    \item Находим её собственные значения $\lambda_1 \ge \lambda_2 \ge \dots \ge 0$.
    \item Сингулярные числа матрицы $A$ равны $\sigma_i = \sqrt{\lambda_i}$. Записываем матрицу $\Sigma$.
    \item Находим ортонормированный базис из собственных векторов $v_1, \dots, v_n$ матрицы $A^T A$. Из них по столбцам составляем матрицу $V$.
    \item Для ненулевых сингулярных чисел вычисляем векторы $u_i = \frac{1}{\sigma_i} A v_i$.
    \item Если нужно, дополняем систему $\{u_i\}$ до ортонормированного базиса всего пространства. Из векторов $u_i$ по столбцам составляем матрицу $U$.
\end{enumerate}


\subsection*{2. Метрические задачи в пространствах многочленов}

В пространстве многочленов $\mathbb{R}[x]_{\le n}$ часто вводится скалярное произведение через определенный интеграл:
\[
(f, g) = \int_a^b f(x)g(x)dx.
\]

\begin{definition}
Расстоянием $d(h, M)$ от элемента $h$ до линейного многообразия $M = y_0 + U$ (где $y_0$ — точка сдвига, $U$ — направляющее подпространство) называется минимум расстояний от $h$ до точек из $M$. Оно вычисляется как длина ортогональной составляющей вектора $(h - y_0)$ относительно подпространства $U$:
\[
d(h, M) = \| (h - y_0) - \text{pr}_U(h - y_0) \|.
\]
\end{definition}

\noindent\textbf{Практический алгоритм поиска расстояния:}
\begin{enumerate}
    \item Выделить направляющее подпространство $U = \langle e_1, \dots, e_k \rangle$ и точку $y_0$.
    \item Перенести начало координат в $y_0$: пусть $v = h - y_0$. Наша задача — найти расстояние от $v$ до $U$.
    \item Составить матрицу Грама $G$ для базиса $U$: $g_{ij} = (e_i, e_j) = \int_a^b e_i(x) e_j(x)dx$.
    \item Составить столбец правых частей $B$: $b_i = (v, e_i) = \int_a^b v(x) e_i(x)dx$.
    \item Решить систему $G \cdot C = B$, где $C$ — столбец коэффициентов проекции.
    \item Найти проекцию: $u = c_1 e_1 + \dots + c_k e_k$.
    \item Найти ортогональную составляющую $w = v - u$ и вычислить искомое расстояние как её длину: $d = \sqrt{(w, w)} = \sqrt{\int_a^b w(x)^2 dx}$.
\end{enumerate}


\subsection*{3. Кривые второго порядка на плоскости}

Общее уравнение линии второго порядка на плоскости имеет вид:
\[
a_{11}x^2 + 2a_{12}xy + a_{22}y^2 + 2a_{13}x + 2a_{23}y + a_{33} = 0.
\]

\noindent\textbf{Ортогональный инвариантный метод приведения к каноническому виду:}
\begin{enumerate}
    \item Выписываем матрицу квадратичной формы старших членов: $A = \begin{pmatrix} a_{11} & a_{12} \\ a_{12} & a_{22} \end{pmatrix}$.
    \item Находим собственные значения $\lambda_1, \lambda_2$ матрицы $A$. Уравнение кривой в новой системе координат примет вид:
    \[
    \lambda_1 (x')^2 + \lambda_2 (y')^2 + 2b_1 x' + 2b_2 y' + a_{33} = 0.
    \]
    \item Находим ортонормированный базис из собственных векторов $v_1, v_2$.
    \item Составляем из них матрицу перехода $C = (v_1 \mid v_2)$. Важно: необходимо выбрать знаки векторов так, чтобы определитель матрицы $C$ был равен $+1$ (преобразование поворота).
    \item Выражаем старые координаты через новые: $\begin{pmatrix} x \\ y \end{pmatrix} = C \begin{pmatrix} x' \\ y' \end{pmatrix}$.
    \item Подставляем $x$ и $y$ в линейную часть исходного уравнения для нахождения $b_1, b_2$.
    \item Выделяем полные квадраты по переменным $x'$ и $y'$, совершая параллельный перенос (сдвиг начала координат), и получаем каноническое уравнение.
\end{enumerate}

\begin{definition}[Эксцентриситет кривой]
Эксцентриситет ($\varepsilon$) — числовая характеристика, определяющая форму кривой:
\begin{itemize}
    \item Для эллипса $\frac{x^2}{a^2} + \frac{y^2}{b^2} = 1$ ($a \ge b$): $\varepsilon = \frac{c}{a} = \sqrt{1 - \frac{b^2}{a^2}}$, где $\varepsilon \in [0, 1)$.
    \item Для гиперболы $\frac{x^2}{a^2} - \frac{y^2}{b^2} = 1$: $\varepsilon = \frac{c}{a} = \sqrt{1 + \frac{b^2}{a^2}}$, где $\varepsilon \in (1, +\infty)$.
    \item Для параболы: $\varepsilon = 1$.
\end{itemize}
\end{definition}


\subsection*{4. Поверхности второго порядка}

\begin{definition}
Поверхность называется \textbf{линейчатой}, если через каждую её точку проходит по крайней мере одна прямая, целиком лежащая на этой поверхности.
\end{definition}

Помимо цилиндров и конусов, существуют две \textit{двояколинейчатые} поверхности (через каждую точку проходят ровно две прямые, целиком лежащие на поверхности):
\begin{enumerate}
    \item \textbf{Однополостный гиперболоид:} $\frac{x^2}{a^2} + \frac{y^2}{b^2} - \frac{z^2}{c^2} = 1$.
    \item \textbf{Гиперболический параболоид («седло»):} $\frac{x^2}{p} - \frac{y^2}{q} = 2z \quad (p>0, q>0)$.
\end{enumerate}

\noindent\textbf{Канонические уравнения основных невырожденных поверхностей:}
\begin{itemize}
    \item Эллипсоид: $\frac{x^2}{a^2} + \frac{y^2}{b^2} + \frac{z^2}{c^2} = 1$
    \item Двуполостный гиперболоид: $\frac{x^2}{a^2} + \frac{y^2}{b^2} - \frac{z^2}{c^2} = -1$
    \item Эллиптический параболоид: $\frac{x^2}{p} + \frac{y^2}{q} = 2z \quad (p>0, q>0)$
    \item Конус второго порядка: $\frac{x^2}{a^2} + \frac{y^2}{b^2} - \frac{z^2}{c^2} = 0$
\end{itemize}
